\documentclass[11pt]{article}

\usepackage[margin=1in]{geometry}
\usepackage{titling}
\usepackage{amsmath,amssymb}
\usepackage{booktabs}
\usepackage{enumitem}
\setlength{\droptitle}{-3em}

\title{HW-6}
\author{
Greyson Knippel \\[4pt]
CPSC 321, Fall 2026\\[4pt]
}
\date{\today}

% attribute / relation names in monospace
\newcommand{\at}[1]{\texttt{#1}}

\begin{document}

\maketitle

%============================================================
\section*{question 1}

\subsection*{(a) relation, primary key, and functional dependencies}

an online game store keeps track of each users game library in one table:
\begin{center}
\at{library(owner\_id, game\_id, username, game\_title, developer, hours\_played)}
\end{center}
Each row says that the user with \at{owner\_id} (whose display name is
\at{username}) owns the game \at{game\_id} (with its \at{game\_title} and
\at{developer}) and has played it for \at{hours\_played} hours.

\medskip
\noindent\textbf{primary key:} \at{(owner\_id, game\_id)}. A user can own a
given game only once, so the pair identifies a row, and
\at{owner\_id, game\_id} $\to$ \at{username, game\_title, developer,
hours\_played}. 

\medskip
\noindent\textbf{non key functional dependencies:}
\begin{enumerate}[label=FD\arabic*:, leftmargin=*]
  \item \at{owner\_id} $\to$ \at{username}. Each account has one display
    name, so every row for the same owner must show the same username.
  \item \at{game\_id} $\to$ \at{game\_title, developer}. Each game in the
    store has one title and one developer, so every row for the same game
    must show the same title and developer.
\end{enumerate}
Both are non trivial and non key dependencies

\subsection*{(b) instance with redundant data}

\begin{center}
\begin{tabular}{llllll}
\toprule
\at{owner\_id} & \at{game\_id} & \at{username} & \at{game\_title} & \at{developer} & \at{hours\_played} \\
\midrule
101 & 413150  & grey     & Stardew Valley & ConcernedApe     & 212 \\
101 & 367520  & grey     & Hollow Knight  & Team Cherry      & 48  \\
202 & 413150  & pixelpat & Stardew Valley & ConcernedApe     & 95  \\
303 & 1145360 & nova     & Hades          & Supergiant Games & 30  \\
\bottomrule
\end{tabular}
\end{center}

\noindent The instance satisfies both FDs: owner 101 is always ``grey'', and
game 413150 is always ``Stardew Valley'' by ConcernedApe. Its also
redundant: ``owner 101 is grey'' is stored twice, and ``413150 is Stardew
Valley by ConcernedApe'' is stored twice. With a real store, a popular game
owned by millions of users would have its title and developer repeated
millions of times.

\subsection*{(c) anomalies}

\begin{description}[leftmargin=1em]
  \item[update anomaly] user 101 changes their username from ``grey'' to
    ``greyson''. If only the first row is updated, the second row still says
    ``grey'', so the table claims owner 101 has two different usernames. This
    breaks \at{owner\_id} $\to$ \at{username}. To make the change correctly we
    have to update every row in user 101s library.

  \item[insertion anomaly] The store adds a new game (say \at{game\_id}
    2000000, ``Starlight Drift'' by a new studio) before anyone has bought it.
    We cannot record its title and developer, because \at{owner\_id} is part
    of the primary key and cannot be NULL, and there is no owner yet. The same
    problem happens for a new user who hasn't bought any games: their
    username can't be stored.

  \item[deletion anomaly] user 303 refunds Hades, so their row is deleted.
    That row was the only one for user 303 and the only one for game
    1145360, so deleting it also loses the fact that user 303 is ``nova''
    and that Hades is made by Supergiant Games, even though we only meant to
    remove one purchase.
\end{description}

%============================================================
\section*{question 2}

\subsection*{(a) decomposition: if $X \to YZ$, then $X \to Y$ and $X \to Z$}

\begin{center}
\begin{tabular}{lll}
\toprule
step & FD & justification \\
\midrule
1 & $X \to YZ$ & given \\
2 & $YZ \to Y$ & reflexivity, since $Y \subseteq YZ$ \\
3 & $X \to Y$  & transitivity on (1) and (2) \\
4 & $YZ \to Z$ & reflexivity, since $Z \subseteq YZ$ \\
5 & $X \to Z$  & transitivity on (1) and (4) \\
\bottomrule
\end{tabular}
\end{center}

\noindent steps (3) and (5) give $X \to Y$ and $X \to Z$ 

\subsection*{(b) pseudotransitivity: if $X \to Y$ and $YW \to Z$, then $XW \to Z$}

\begin{center}
\begin{tabular}{lll}
\toprule
step & FD & justification \\
\midrule
1 & $X \to Y$   & given \\
2 & $XW \to YW$ & augmentation of (1) with $W$ \\
3 & $YW \to Z$  & given \\
4 & $XW \to Z$  & transitivity on (2) and (3) \\
\bottomrule
\end{tabular}
\end{center}

\noindent step (4) gives $XW \to Z$

%============================================================
\section*{question 3}

\noindent\textbf{method.} for a set of attributes $X$, write $X^+$ for its
closure: every attribute that $X$ determines. to find it, start with $X$
(reflexivity) and keep applying FDs: if some FD $Y \to Z$ has $Y$ already in
the closure, add $Z$ (this is augmentation/decomposition plus transitivity).
$X$ is a superkey exactly when $X^+$ contains every attribute of $R$, and a
candidate key when no proper subset of $X$ is also a superkey.

\medskip
\noindent a useful shortcut: an attribute that never appears on the
right-hand side of any FD can't be determined by anything else, so it has to
be in \emph{every} key.

%------------------------------------------------------------
\subsection*{(a) $R(a,b,c,d)$, FDs $\{b \to c,\ d \to a\}$}

\noindent\textbf{candidate key: $bd$ (the only one).}

\begin{itemize}
  \item \textbf{determines everything:} $\{b,d\}$
    $\xrightarrow{b \to c}$ $\{b,c,d\}$
    $\xrightarrow{d \to a}$ $\{a,b,c,d\}$. so $bd^+$ is all of $R$.
  \item \textbf{minimal:} $b^+ = \{b,c\}$ (missing $a, d$) and
    $d^+ = \{a,d\}$ (missing $b, c$), so neither $b$ nor $d$ alone is a
    superkey.
  \item \textbf{no other keys:} $b$ and $d$ never appear on the right-hand
    side of an FD, so every superkey must contain both of them, i.e.\ every
    superkey contains $bd$. any superkey other than $bd$ is then a proper
    superset of $bd$ and isn't minimal.
\end{itemize}

%------------------------------------------------------------
\subsection*{(b) $R(a,b,c,d,e)$, FDs $\{a \to bc,\ cd \to e,\ b \to d,\ e \to a\}$}

\noindent\textbf{candidate keys: $a$, $e$, $bc$, $cd$.}

\medskip
\noindent here every attribute shows up on some right-hand side, so the
shortcut doesn't help and we check closures directly.

\begin{itemize}
  \item \textbf{$a$:} $\{a\}$
    $\xrightarrow{a \to bc}$ $\{a,b,c\}$
    $\xrightarrow{b \to d}$ $\{a,b,c,d\}$
    $\xrightarrow{cd \to e}$ $\{a,b,c,d,e\}$.
    a single attribute is automatically minimal (its only proper subset is
    $\emptyset$, which determines nothing).
  \item \textbf{$e$:} $\{e\}$
    $\xrightarrow{e \to a}$ $\{a,e\}$, and from $a$ we get everything as
    above, so $e^+ = \{a,b,c,d,e\}$. minimal for the same reason.
  \item \textbf{$bc$:} $\{b,c\}$
    $\xrightarrow{b \to d}$ $\{b,c,d\}$
    $\xrightarrow{cd \to e}$ $\{b,c,d,e\}$
    $\xrightarrow{e \to a}$ $\{a,b,c,d,e\}$.
    minimal since $b^+ = \{b,d\}$ and $c^+ = \{c\}$.
  \item \textbf{$cd$:} $\{c,d\}$
    $\xrightarrow{cd \to e}$ $\{c,d,e\}$
    $\xrightarrow{e \to a}$ $\{a,c,d,e\}$
    $\xrightarrow{a \to bc}$ $\{a,b,c,d,e\}$.
    minimal since $c^+ = \{c\}$ and $d^+ = \{d\}$ (no FD has just $c$ or
    just $d$ on the left).
\end{itemize}

\noindent\textbf{no other keys.} take any superkey $K$.
\begin{itemize}
  \item if $K$ contains $a$ or $e$, then $K$ contains the key $a$ or the key
    $e$, so it's only minimal if $K = a$ or $K = e$.
  \item otherwise $K \subseteq \{b,c,d\}$. the singletons fail
    ($b^+ = \{b,d\}$, $c^+ = \{c\}$, $d^+ = \{d\}$). of the pairs, $bc$ and
    $cd$ are keys, and $bd^+ = \{b,d\}$ is not a superkey. the only triple,
    $bcd$, contains $bc$, so it isn't minimal.
\end{itemize}
so the candidate keys are exactly $a$, $e$, $bc$, and $cd$.

%------------------------------------------------------------
\subsection*{(c) $R(a,b,c,d,e)$, FDs $\{c \to b,\ bd \to e,\ a \to d,\ e \to a\}$}

\noindent\textbf{candidate keys: $ac$, $cd$, $ce$.}

\medskip
\noindent $c$ never appears on a right-hand side, so $c$ is in every key. on
its own, $c^+ = \{b,c\}$, which is not everything, so we need to add more.

\begin{itemize}
  \item \textbf{$ac$:} $\{a,c\}$
    $\xrightarrow{c \to b}$ $\{a,b,c\}$
    $\xrightarrow{a \to d}$ $\{a,b,c,d\}$
    $\xrightarrow{bd \to e}$ $\{a,b,c,d,e\}$.
    minimal since $a^+ = \{a,d\}$ and $c^+ = \{b,c\}$.
  \item \textbf{$cd$:} $\{c,d\}$
    $\xrightarrow{c \to b}$ $\{b,c,d\}$
    $\xrightarrow{bd \to e}$ $\{b,c,d,e\}$
    $\xrightarrow{e \to a}$ $\{a,b,c,d,e\}$.
    minimal since $c^+ = \{b,c\}$ and $d^+ = \{d\}$.
  \item \textbf{$ce$:} $\{c,e\}$
    $\xrightarrow{c \to b}$ $\{b,c,e\}$
    $\xrightarrow{e \to a}$ $\{a,b,c,e\}$
    $\xrightarrow{a \to d}$ $\{a,b,c,d,e\}$.
    minimal since $c^+ = \{b,c\}$ and $e^+ = \{a,d,e\}$.
\end{itemize}

\noindent\textbf{no other keys.} every key contains $c$. if a superkey $K$
also contains $a$, $d$, or $e$, then it contains one of the keys $ac$, $cd$,
or $ce$, so it's only minimal if it equals that key. otherwise
$K \subseteq \{b,c\}$, but $c^+ = bc^+ = \{b,c\}$, so neither $c$ nor $bc$ is
a superkey. so the candidate keys are exactly $ac$, $cd$, and $ce$.

\end{document}
